\chapter{Supramolecular Chemistry}\label{ch:b3:supramolecular}

In 1962 Charles Pedersen, looking for something else, isolated a cyclic
polyether that dissolved potassium permanganate in benzene and turned it
purple. The ring had wrapped itself around the potassium ion and hidden its
charge behind a hydrocarbon surface, dragging the permanganate along as the
counter-ion. Chemistry beyond the molecule, the chemistry of assemblies held
together by forces weaker than covalent bonds, began with that purple solution.
This chapter measures those forces, explains why some \omterm{def:b3:supramolecular:supramolecular}{hosts} choose their \omterm{def:b3:supramolecular:supramolecular}{guests}
so well, shows how binding constants are measured, and ends with molecules
threaded and interlocked into machines.

\begin{recall}
The Year 1 volume treated intermolecular forces, the hydrogen bond, solvation,
hydrophobic molecules, complexes, formation constants and chelates; the Year 2
volume the chemical potential, Gibbs energies and least-squares fitting.
\Cref{ch:b3:advanced-nmr} described fast and slow exchange and coalescence,
\cref{ch:b3:rate-theories} the \omterm{thm:b3:rate-theories:eyring}{Eyring equation}, and
\cref{ch:b3:partition-functions} the translational entropy of a gas.
\end{recall}

\section{Non-covalent interactions}

\begin{definition}[Supramolecular chemistry]\label{def:b3:supramolecular:supramolecular}
\emph{Supramolecular chemistry}\index{supramolecular chemistry} is the chemistry
of assemblies of molecules held together by non-covalent interactions. In a
\emph{host--guest complex}\index{host--guest complex}, the \emph{host}\index{host}
is the larger molecule, with a cavity or a cleft and binding sites converging on
it, and the \emph{guest}\index{guest} the smaller molecule or ion it binds.
\end{definition}

The interactions are those of the Year 1 volume, ion--dipole, dipole--dipole,
hydrogen bonds and London forces, plus a few with names of their own.

\begin{definition}[Non-covalent interactions]\label{def:b3:supramolecular:interactions}
\emph{$\pi$ stacking}\index{$\pi$ stacking} is the attraction between flat aromatic
rings lying face to face, usually offset or tilted. A \emph{cation--$\pi$
interaction}\index{cation--$\pi$ interaction} is the attraction between a cation and
the electron-rich face of an aromatic ring. A \emph{halogen bond}\index{halogen bond}
is the attraction between the electron-poor tip of a covalently bound
halogen atom, opposite its bond, and a Lewis base. The \emph{hydrophobic
effect}\index{hydrophobic effect} is the tendency of non-polar surfaces to come
together in water, driven mostly by the release of the ordered water molecules
around them.
\end{definition}

The ion--dipole and cation--$\pi$ interactions are the strongest, and depend
heavily on the solvent that competes for the ion; hydrogen bonds come next,
then $\pi$ stacking and \omterm{def:b3:supramolecular:interactions}{halogen bonds}, then London forces, weak one by one but
summed over large contact surfaces. In water, the \omterm{def:b3:supramolecular:interactions}{hydrophobic effect} often
dominates: a \omterm{def:b3:supramolecular:supramolecular}{host} built to bind a \omterm{def:b3:supramolecular:supramolecular}{guest} by hydrogen bonds in chloroform may not
bind it at all in water, whose own hydrogen bonds compete.

\section{Molecular recognition}

\begin{definition}[Molecular recognition]\label{def:b3:supramolecular:recognition}
\emph{Molecular recognition}\index{molecular recognition} is the selective
binding of a \omterm{def:b3:supramolecular:supramolecular}{guest} by a \omterm{def:b3:supramolecular:supramolecular}{host}. It relies on \emph{complementarity}\index{complementarity},
a match of size, shape and binding sites between \omterm{def:b3:supramolecular:supramolecular}{host} and \omterm{def:b3:supramolecular:supramolecular}{guest}, and is
strengthened by \emph{preorganisation}\index{preorganisation}, a \omterm{def:b3:supramolecular:supramolecular}{host} whose
binding sites are already held in the geometry of the complex before binding.
\end{definition}

The association constant $K_a = [\mathrm{HG}]/([\mathrm H][\mathrm G])$ is a formation constant in the sense of
the Year 1 volume, and $\Delta G^\circ = -RT\ln K_a$. A flexible \omterm{def:b3:supramolecular:supramolecular}{host} must freeze its
conformation on binding, which costs entropy; a rigid, preorganised one has
paid that price in its synthesis.

\begin{definition}[Chelate and macrocyclic effects]\label{def:b3:supramolecular:macrocyclic}
The \emph{chelate effect}\index{chelate effect} is the greater stability of a
complex with a polydentate ligand than with the equivalent number of
monodentate ligands with the same donor atoms. The \emph{macrocyclic
effect}\index{macrocyclic effect} is the further stabilisation when the
polydentate ligand is a ring rather than an open chain.
\end{definition}

\begin{proposition}[The chelate effect is mainly entropic]\label{prop:b3:supramolecular:chelate-entropy}
In $\mathrm{[M(NH_3)_6]^{2+} + 3\,en \rightleftharpoons [M(en)_3]^{2+} + 6\,NH_3}$, where en is ethane-1,2-diamine, the same six
metal--nitrogen bonds are present on both sides, and the reaction is driven by
the increase in the number of free molecules.
\end{proposition}

\begin{proof}[Argued]
The bonds made and broken are of the same kind, so $\Delta_rH^\circ$ is small. The number of
independent molecules goes from four to seven: three more translational
degrees of freedom of a whole molecule, each worth tens of joules per kelvin
and per mole in solution (\cref{ch:b3:partition-functions}), so $\Delta_rS^\circ > 0$ and $K > 1$.
Seen kinetically, once one end of a chelate is bound, the other is held near the
metal at a high effective concentration, and binds at once if it lets go.
\end{proof}

\begin{proposition}[Enthalpy--entropy compensation]\label{prop:b3:supramolecular:compensation}
In a series of related \omterm{def:b3:supramolecular:supramolecular}{host--guest complexes}, a more negative $\Delta H^\circ$ of binding
usually comes with a more negative $\Delta S^\circ$, so that $\Delta G^\circ$ varies less than either (an
experimental observation).
\end{proposition}

A tighter complex binds more strongly but restricts the motions of \omterm{def:b3:supramolecular:supramolecular}{host} and
\omterm{def:b3:supramolecular:supramolecular}{guest} more; a stronger hydrogen bond to the \omterm{def:b3:supramolecular:supramolecular}{guest} often means a weaker one to
the water it displaced. The Gibbs energy, which sets selectivity, is the small
difference of two large terms.

\begin{definition}[Hosts]\label{def:b3:supramolecular:hosts}
A \emph{crown ether}\index{crown ether} is a macrocyclic polyether, such as
18-crown-6, $\ce{(CH2CH2O)6}$. A \emph{cryptand}\index{cryptand} is a bicyclic
polyether with two nitrogen bridgeheads, which encloses a cation in three
dimensions in a \emph{cryptate}\index{cryptate}. A
\emph{cyclodextrin}\index{cyclodextrin} is a cyclic oligomer of glucose units, in
the shape of a truncated cone with a hydrophobic cavity and hydroxy groups on
its rims. A \emph{calixarene}\index{calixarene} is a cup-shaped macrocycle of
phenol units joined by methylene bridges.
\end{definition}

\begin{omfigure}
\begin{tikzpicture}[font=\small]
  % 18-crown-6 with K+
  \foreach \k in {0,...,17} {
    \pgfmathsetmacro\ang{90 + 20*\k}
    \pgfmathsetmacro\nxt{110 + 20*\k}
    \draw[black!70] (\ang:1.55) -- (\nxt:1.55);
  }
  \foreach \k in {0,3,...,15} {
    \pgfmathsetmacro\ang{90 + 20*\k}
    \node[fill=white, inner sep=0.5pt, text=omThm] (o\k) at (\ang:1.55) {O};
    \draw[omThm!60, densely dashed] (0,0) -- (\ang:1.32);
  }
  \node[circle, fill=omDef!25, draw=omDef, inner sep=1pt] at (0,0) {$\ce{K+}$};
  \node[font=\scriptsize] at (0,-2.1) {18-crown-6 with $\ce{K+}$};
  % [2.2.2]cryptand with K+
  \begin{scope}[xshift=5.2cm]
    \node (nt) at (0,1.7) {N};
    \node (nb) at (0,-1.7) {N};
    \foreach \x/\n in {-1.25/a, 0/b, 1.25/c} {
      \node[text=omThm, fill=white, inner sep=0.5pt] (oa\n) at (\x,0.55) {O};
      \node[text=omThm, fill=white, inner sep=0.5pt] (ob\n) at (\x,-0.55) {O};
      \draw[black!70] (nt) to[out=-90+40*\x, in=90] (oa\n) -- (ob\n) to[out=-90, in=90-40*\x] (nb);
    }
    \node[circle, fill=omDef!25, draw=omDef, inner sep=1pt] at (0.62,0) {$\ce{K+}$};
    \node[font=\scriptsize] at (0,-2.4) {[2.2.2]cryptand: $\ce{K+}$ cryptate};
  \end{scope}
\end{tikzpicture}

{\small Left: 18-crown-6, a ring of twelve carbon atoms (corners) and six oxygen
atoms, with a potassium ion at its centre bound by the six oxygens. Right: the
[2.2.2]\omterm{def:b3:supramolecular:hosts}{cryptand}, three chains of two ether oxygens between two nitrogen
bridgeheads, schematic: the cation is held inside a three-dimensional cage,
by six oxygens and the two nitrogens.}
\end{omfigure}

\omterm{def:b3:supramolecular:hosts}{Crown ethers} select cations by size: 18-crown-6 binds $\ce{K+}$ more strongly than the
smaller $\ce{Na+}$, which cannot reach all six oxygens at once, or the larger $\ce{Cs+}$, which
sits above the ring. The \omterm{def:b3:supramolecular:hosts}{cryptands}, more preorganised and wrapping the ion
completely, bind more strongly still and select more sharply. Hidden inside
the \omterm{def:b3:supramolecular:supramolecular}{host}, the cation drags its anion into non-polar solvents, where the anion,
unsolvated, becomes very reactive: the basis of phase-transfer catalysis.

\section{Measuring binding}

\begin{proposition}[Fast exchange]\label{prop:b3:supramolecular:fast-exchange}
When a \omterm{def:b3:supramolecular:supramolecular}{host} exchanges between its free and bound states much faster than the
difference of their resonance frequencies, its NMR signal appears at the average
$\delta_{\mathrm{obs}} = x_{\mathrm{free}}\delta_{\mathrm{free}} + x_{\mathrm{bound}}\delta_{\mathrm{bound}}$, where $x$ are the fractions of \omterm{def:b3:supramolecular:supramolecular}{host} in each state.
\end{proposition}

\begin{proof}
Each \omterm{def:b3:supramolecular:supramolecular}{host} molecule visits both states many times during the measurement; its
precession frequency, averaged over time, is the weighted mean of the two, with
weights the fractions of time spent in each, which at equilibrium are the
fractions of molecules in each state (\cref{ch:b3:advanced-nmr}).
\end{proof}

\begin{theorem}[Exact 1:1 isotherm]\label{thm:b3:supramolecular:isotherm}
For $\mathrm{H + G \rightleftharpoons HG}$ with total concentrations $H_0$ and $G_0$ and association constant $K$,
\[
  [\mathrm{HG}] = \frac{1}{2}\Bigl[\Bigl(H_0 + G_0 + \frac1K\Bigr) - \sqrt{\Bigl(H_0 + G_0 + \frac1K\Bigr)^2 - 4H_0G_0}\,\Bigr].
\]
\end{theorem}

\begin{proof}
With $c = [\mathrm{HG}]$, $K = c/((H_0 - c)(G_0 - c))$, that is $c^2 - (H_0 + G_0 + 1/K)c + H_0G_0 = 0$. Of the two roots,
only the smaller is below both $H_0$ and $G_0$ (their product is $H_0G_0$ and their sum
exceeds $H_0 + G_0$): it is the one with the minus sign.
\end{proof}

\begin{omfigure}
\begin{tikzpicture}
\begin{axis}[name=a, width=0.48\linewidth, height=5.5cm, xmin=0, xmax=6, ymin=0, ymax=1.05,
  axis lines=left, xlabel={equivalents of guest, $G_0/H_0$}, ylabel={fraction of host bound},
  tick label style={font=\scriptsize}, label style={font=\small},
  ]
\addplot[omMeth, thick] table[x=equiv, y=K4] {figdata/out/bachelor-3-binding-titration.dat};
\addplot[omDef, thick] table[x=equiv, y=K3] {figdata/out/bachelor-3-binding-titration.dat};
\addplot[omThm, thick] table[x=equiv, y=K2] {figdata/out/bachelor-3-binding-titration.dat};
\node[font=\scriptsize, omMeth, anchor=north] at (axis cs:4,0.94) {$K = 10^4$};
\node[font=\scriptsize, omDef, anchor=north] at (axis cs:4,0.68) {$K = 10^3$};
\node[font=\scriptsize, omThm, anchor=north] at (axis cs:4,0.27) {$K = 10^2$};
\end{axis}
\begin{axis}[at={(a.east)}, anchor=west, xshift=1.4cm, width=0.48\linewidth, height=5.5cm,
  xmin=0, xmax=1, ymin=0, ymax=1.05, axis lines=left,
  xlabel={mole fraction of host}, ylabel={$[\text{complex}]/c_{\text{total}}$},
  tick label style={font=\scriptsize}, label style={font=\small}]
\addplot[omDef, thick] table[x=x, y=HG] {figdata/out/bachelor-3-binding-titration-b.dat};
\addplot[omThm, thick] table[x=x, y=HG2] {figdata/out/bachelor-3-binding-titration-b.dat};
\draw[black!35, dashed] (axis cs:0.5,0) -- (axis cs:0.5,0.55) (axis cs:0.3333,0) -- (axis cs:0.3333,0.38);
\node[font=\scriptsize, omDef, anchor=south] at (axis cs:0.5,0.52) {HG};
\node[font=\scriptsize, omThm, anchor=south] at (axis cs:0.25,0.36) {$\mathrm{HG_2}$};
\end{axis}
\end{tikzpicture}

{\small Left: fraction of a \omterm{def:b3:supramolecular:supramolecular}{host} bound during a titration at $H_0 = \qty{1.0}{mmol/L}$ for three
association constants (in $\unit{L.mol^{-1}}$): the larger $K$, the sharper the break at one equivalent;
when $KH_0 \gg 1$ the curve only says that $K$ is large. Right: \omterm{def:b3:supramolecular:job}{Job plots}, the
concentration of complex at a constant total concentration of \omterm{def:b3:supramolecular:supramolecular}{host} plus \omterm{def:b3:supramolecular:supramolecular}{guest},
peaking at a \omterm{def:b3:supramolecular:supramolecular}{host} fraction of 1/2 for a 1:1 and 1/3 for a 1:2 complex (model
constants).}
\end{omfigure}

\begin{method}[Binding constant from an NMR titration]\label{met:b3:supramolecular:nmr-titration}
\begin{enumerate}
  \item Keep the \omterm{def:b3:supramolecular:supramolecular}{host} concentration $H_0$ fixed and add \omterm{def:b3:supramolecular:supramolecular}{guest}, from 0 to several
        equivalents; record a \omterm{def:b3:supramolecular:supramolecular}{host} signal at each point (fast exchange).
  \item Model: $\delta = \delta_{\mathrm{free}} + (\delta_{\mathrm{bound}} - \delta_{\mathrm{free}})[\mathrm{HG}]/H_0$, with $[\mathrm{HG}]$ from the exact isotherm.
  \item Fit $K$ and $\delta_{\mathrm{bound}}$ by nonlinear least squares, minimising the sum of the squared
        residuals; take the standard errors from the curvature of that sum
        (the covariance matrix).
  \item Check the residuals for a trend (wrong stoichiometry) and that the
        bound fraction spans roughly 20 to 80\,\% (the data then constrain $K$).
\end{enumerate}
\end{method}

\begin{definition}[Job plot]\label{def:b3:supramolecular:job}
A \emph{Job plot}\index{Job plot} (method of continuous variations) is a plot of
the concentration of a complex, or of a signal proportional to it, against
the mole fraction of one component, at a constant total concentration of the
two.
\end{definition}

\begin{proposition}[Maximum of a Job plot]\label{prop:b3:supramolecular:job}
For a single complex $\mathrm{H_nG_m}$, the \omterm{def:b3:supramolecular:job}{Job plot} is maximal at the \omterm{def:b3:supramolecular:supramolecular}{host} mole fraction
$x = n/(n + m)$, whatever the binding constant.
\end{proposition}

\begin{proof}
With total concentration $T$, $H_t = xT$, $G_t = (1 - x)T$, complex $c$, free $[\mathrm H] = H_t - nc$,
$[\mathrm G] = G_t - mc$, and $c = \beta[\mathrm H]^n[\mathrm G]^m$. Differentiating with respect to $x$ at the maximum, where
$\dd c/\dd x = 0$: $0 = \beta\bigl(n[\mathrm H]^{n-1}[\mathrm G]^m T - m[\mathrm H]^n[\mathrm G]^{m-1}T\bigr)$, so $n[\mathrm G] = m[\mathrm H]$. Substituting,
$n(G_t - mc) = m(H_t - nc)$, so $nG_t = mH_t$, that is $n(1 - x) = mx$: $x = n/(n + m)$.
\end{proof}

\begin{method}[Running a Job plot]\label{met:b3:supramolecular:job}
\begin{enumerate}
  \item Prepare equimolar stock solutions of \omterm{def:b3:supramolecular:supramolecular}{host} and \omterm{def:b3:supramolecular:supramolecular}{guest}.
  \item Mix them in proportions $x$ from 0 to 1 at a constant total volume.
  \item Measure a property proportional to the complex (an absorbance corrected
        for the free species, or $x\,\Delta\delta$ in NMR).
  \item Read the stoichiometry from the position of the maximum.
\end{enumerate}
\end{method}

\begin{definition}[Isothermal titration calorimetry]\label{def:b3:supramolecular:itc}
\emph{Isothermal titration calorimetry}\index{isothermal titration calorimetry}
(ITC) measures the heat released or absorbed at each injection of a \omterm{def:b3:supramolecular:supramolecular}{guest}
solution into a \omterm{def:b3:supramolecular:supramolecular}{host} solution held at constant temperature.
\end{definition}

\begin{proposition}[Heat per injection]\label{prop:b3:supramolecular:itc}
Neglecting dilution, the heat of the $i$-th injection into a cell of volume $V_0$ is
$q_i = \Delta H^\circ V_0\bigl([\mathrm{HG}]_i - [\mathrm{HG}]_{i-1}\bigr)$.
\end{proposition}

\begin{proof}
The heat is the reaction enthalpy times the amount of complex formed during the
injection, and that amount is the change of $[\mathrm{HG}]$ times the cell volume.
\end{proof}

\begin{method}[Reading an ITC isotherm]\label{met:b3:supramolecular:itc}
\begin{enumerate}
  \item Plot the heat per mole of injected \omterm{def:b3:supramolecular:supramolecular}{guest} against the molar ratio.
  \item The height of the plateau before saturation gives $\Delta H^\circ$; the molar ratio at
        the inflection gives the stoichiometry; the steepness gives $K$.
  \item Fit all three with the exact isotherm; then $\Delta G^\circ = -RT\ln K$ and
        $T\Delta S^\circ = \Delta H^\circ - \Delta G^\circ$.
\end{enumerate}
\end{method}

A single ITC experiment thus gives the whole thermodynamic signature of the
binding, enthalpy and entropy, which an NMR or UV--visible titration gives only
through a series of temperatures.

\begin{inthelab}[An NMR titration]
A solution of the \omterm{def:b3:supramolecular:supramolecular}{host} in a deuterated solvent, about a millimole per litre,
is placed in an NMR tube and its spectrum recorded. A concentrated solution of
the \omterm{def:b3:supramolecular:supramolecular}{guest} made up in the \omterm{def:b3:supramolecular:supramolecular}{host} solution, so that the \omterm{def:b3:supramolecular:supramolecular}{host} concentration stays
constant, is added in small portions by microsyringe; after each addition the
tube is shaken, equilibrated at the probe temperature, and the spectrum
recorded. The shift of a \omterm{def:b3:supramolecular:supramolecular}{host} signal that moves is read at each point and
fitted.
\end{inthelab}

\section{Hosts}

\omterm{def:b3:supramolecular:hosts}{Cyclodextrins}, made by \omterm{def:b3:complex-kinetics:enzyme}{enzymes} from starch, have six, seven or eight glucose
units ($\alpha$, $\beta$, $\gamma$). Their hydrophobic cavity holds non-polar parts of \omterm{def:b3:supramolecular:supramolecular}{guests} in
water: they solubilise drugs, carry fragrances and remove cholesterol from
cell membranes. \omterm{def:b3:supramolecular:hosts}{Calixarenes} and the pumpkin-shaped cucurbiturils, rigid
macrocycles of glycoluril units with carbonyl-lined portals, bind cations and
organic cations very strongly in water.

\begin{omfigure}
\begin{tikzpicture}[font=\small]
  \draw[black!70, fill=omDef!8] (-1.9,1.2) -- (-1.2,-1.2) arc[start angle=180, end angle=360, x radius=1.2, y radius=0.3]
    -- (1.9,1.2);
  \draw[black!70, fill=omDef!18] (0,1.2) ellipse[x radius=1.9, y radius=0.45];
  \draw[black!40, dashed] (-1.2,-1.2) arc[start angle=180, end angle=0, x radius=1.2, y radius=0.3];
  \node[draw=omThm, thick, regular polygon, regular polygon sides=6, minimum size=8mm, inner sep=0pt] at (0,0.2) {};
  \draw[omThm, thick] (0,0.62) -- (0,2.0);
  \node[font=\scriptsize, omThm, anchor=west] at (0.1,2.0) {guest};
  \node[font=\scriptsize, anchor=west] at (2.0,1.2) {wide rim: secondary OH};
  \node[font=\scriptsize, anchor=west] at (1.3,-1.25) {narrow rim: primary OH};
  \node[font=\scriptsize, anchor=east] at (-1.75,0) {hydrophobic cavity};
\end{tikzpicture}

{\small A \omterm{def:b3:supramolecular:hosts}{cyclodextrin}, drawn as the truncated cone it forms, with an aromatic
\omterm{def:b3:supramolecular:supramolecular}{guest} in its cavity. The hydroxy groups on the two rims make the outside
water-soluble; the inside, lined with C--H groups and glycosidic oxygens, is
non-polar.}
\end{omfigure}

\section{Self-assembly and molecular machines}

\begin{definition}[Self-assembly]\label{def:b3:supramolecular:self-assembly}
\emph{Self-assembly}\index{self-assembly} is the spontaneous and reversible
organisation of components into a defined structure, under the control of the
information contained in the components themselves: their shapes and binding
sites.
\end{definition}

Metal ions with fixed coordination angles and rigid bridging ligands assemble
into squares, cages and capsules in a single step, because each bond can open
and close until the most stable structure, without strain or dangling sites,
is reached. Base pairing in DNA, and the folding of proteins, are the same idea
on a grander scale.

\begin{definition}[Mechanically interlocked molecules]\label{def:b3:supramolecular:interlocked}
A \emph{mechanically interlocked molecule}\index{mechanically interlocked molecule}
has parts that are not bonded to each other but cannot be separated
without breaking a covalent bond. A \emph{rotaxane}\index{rotaxane} is a ring
threaded on an axle with bulky stoppers at both ends; a
\emph{catenane}\index{catenane} consists of interlocked rings.
\end{definition}

\begin{definition}[Template synthesis]\label{def:b3:supramolecular:template}
A \emph{template synthesis}\index{template synthesis} uses a temporary
interaction, often with a metal ion, to hold reactive components in the
geometry needed for a reaction that would otherwise be improbable, such as the
closure of a ring around another ring.
\end{definition}

\begin{omfigure}
\begin{tikzpicture}[font=\small, >=Stealth]
  % step 1: two crescents around Cu+
  \draw[omDef, line width=2pt] (0.3,0) arc[start angle=0, end angle=300, radius=0.75];
  \draw[omThm, line width=2pt] (-0.3,0) arc[start angle=180, end angle=-120, radius=0.75];
  \fill[omProp] (0,0) circle (0.13);
  \node[font=\scriptsize] at (0,-1.35) {$\ce{Cu+}$ holds two};
  \node[font=\scriptsize] at (0,-1.7) {open ligands crossed};
  \draw[->, thick] (1.4,0) -- (2.3,0) node[midway, above, font=\scriptsize] {close};
  % step 2: closed interlocked rings with Cu
  \begin{scope}[xshift=3.6cm]
    \draw[omDef, line width=2pt] (-0.35,0) circle (0.75);
    \draw[omThm, line width=2pt] (0.35,0) circle (0.75);
    \draw[white, line width=5pt] ([shift={(52:0.75)}]-0.35,0) arc[start angle=52, end angle=72, radius=0.75];
    \draw[omDef, line width=2pt] ([shift={(48:0.75)}]-0.35,0) arc[start angle=48, end angle=76, radius=0.75];
    \fill[omProp] (0,0) circle (0.13);
    \node[font=\scriptsize] at (0,-1.35) {catenate};
  \end{scope}
  \draw[->, thick] (5.0,0) -- (5.9,0) node[midway, above, font=\scriptsize] {$-\ce{Cu+}$};
  \begin{scope}[xshift=7.2cm]
    \draw[omDef, line width=2pt] (-0.35,0) circle (0.75);
    \draw[omThm, line width=2pt] (0.35,0) circle (0.75);
    \draw[white, line width=5pt] ([shift={(52:0.75)}]-0.35,0) arc[start angle=52, end angle=72, radius=0.75];
    \draw[omDef, line width=2pt] ([shift={(48:0.75)}]-0.35,0) arc[start angle=48, end angle=76, radius=0.75];
    \node[font=\scriptsize] at (0,-1.35) {[2]catenane};
  \end{scope}
  % rotaxane
  \begin{scope}[xshift=10.2cm]
    \draw[black!70, line width=1.6pt] (-1.1,0) -- (1.1,0);
    \fill[black!60] (-1.25,0) circle (0.2) (1.25,0) circle (0.2);
    \draw[omThm, line width=2pt] (0,0) ellipse[x radius=0.18, y radius=0.55];
    \node[font=\scriptsize] at (0,-1.35) {[2]rotaxane};
  \end{scope}
\end{tikzpicture}

{\small \omterm{def:b3:supramolecular:template}{Template synthesis} of a \omterm{def:b3:supramolecular:interlocked}{catenane} (schematic): a copper(I) ion holds two
phenanthroline-based ligands crossed at right angles in its tetrahedral
coordination; closing each into a ring gives two interlocked rings around the
metal, and removing the copper leaves the free \omterm{def:b3:supramolecular:interlocked}{catenane}. Right: a \omterm{def:b3:supramolecular:interlocked}{rotaxane}, a
ring trapped on an axle by two stoppers.}
\end{omfigure}

\begin{definition}[Molecular machines]\label{def:b3:supramolecular:machine}
A \emph{molecular machine}\index{molecular machine} is an assembly whose
components move relative to one another in a controlled way in response to a
stimulus (light, a redox change, a change of pH), and can be made to do work.
\end{definition}

In a molecular shuttle, a ring on a \omterm{def:b3:supramolecular:interlocked}{rotaxane} axle sits at one of two stations; a
stimulus that weakens its binding to the first sends it to the second, and the
reverse stimulus sends it back. Light-driven rotary motors based on
overcrowded alkenes turn in one direction only, by alternating photochemical
isomerisations and thermal steps that cannot run backwards. The rate of
shuttling between stations is measured by variable-temperature NMR, from the
coalescence of the signals of the two forms (\cref{ch:b3:advanced-nmr}).

\begin{safety}{\ghs[0.62cm]{GHS07}}
% ledger: ghs4:18C6
18-Crown-6: harmful if swallowed, irritating to the skin, eyes and airways.
\omterm{def:b3:supramolecular:hosts}{Crown ethers} carry metal salts into organic solvents and through skin; they are
weighed and handled with gloves.
\end{safety}

\begin{history}[Supramolecular chemistry and its machines]
Charles Pedersen's \omterm{def:b3:supramolecular:hosts}{crown ethers} (1967), Jean-Marie Lehn's \omterm{def:b3:supramolecular:hosts}{cryptands} and Donald
Cram's preorganised \omterm{def:b3:supramolecular:supramolecular}{hosts} founded the field; the three shared the 1987 Nobel
Prize in Chemistry. Jean-Pierre Sauvage made \omterm{def:b3:supramolecular:interlocked}{catenanes} in high yield by the
copper template in 1983, Fraser Stoddart built \omterm{def:b3:supramolecular:interlocked}{rotaxane} shuttles, and Ben
Feringa a light-driven rotary motor in 1999; they shared the 2016 Nobel Prize in
Chemistry for the design and synthesis of \omterm{def:b3:supramolecular:machine}{molecular machines}.
\end{history}

\section{Exercises}

\begin{exercise}[$\star$]\label{exo:b3:supramolecular:1}
Name the main interaction in: $\ce{K+}$ in 18-crown-6; two benzene rings face to face;
$\ce{K+}$ above a benzene ring; iodopentafluorobenzene with pyridine; a guanine--cytosine
pair.
\end{exercise}

\begin{exercise}[$\star$]\label{exo:b3:supramolecular:2}
Compute $\Delta G^\circ$ at \qty{298}{K} for a \omterm{def:b3:supramolecular:supramolecular}{host--guest complex} with $K_a = \qty{1.0e4}{L.mol^{-1}}$.
\end{exercise}

\begin{exercise}[$\star$]\label{exo:b3:supramolecular:3}
In methanol, $\log K = 6.1$ for $\ce{K+}$ and 4.4 for $\ce{Na+}$ with 18-crown-6 (data of the exercise).
Compute the selectivity and the difference in $\Delta G^\circ$.
\end{exercise}

\begin{exercise}[$\star$]\label{exo:b3:supramolecular:4}
A \omterm{def:b3:supramolecular:job}{Job plot} peaks at a \omterm{def:b3:supramolecular:supramolecular}{host} mole fraction of 0.33. What is the stoichiometry?
\end{exercise}

\begin{exercise}[$\star\star$]\label{exo:b3:supramolecular:5}
A \omterm{def:b3:supramolecular:supramolecular}{host} at \qty{2.0}{mmol/L} with $K = \qty{1500}{L.mol^{-1}}$, $\delta_{\mathrm{free}} = \qty{3.560}{ppm}$ and $\delta_{\mathrm{bound}} = \qty{3.720}{ppm}$ receives one
equivalent of \omterm{def:b3:supramolecular:supramolecular}{guest}. Compute the bound fraction and the observed shift.
\end{exercise}

\begin{exercise}[$\star\star$]\label{exo:b3:supramolecular:6}
For nickel(II), $\log\beta_6 = 8.6$ with ammonia and $\log\beta_3 = 18.3$ with ethane-1,2-diamine (data of the
exercise). Compute the constant and $\Delta G^\circ$ at \qty{298}{K} of
$\ce{[Ni(NH3)6]^2+} + 3\,\mathrm{en} \rightleftharpoons \ce{[Ni(en)3]^2+} + \ce{6NH3}$, and say what drives it.
\end{exercise}

\begin{exercise}[$\star\star$]\label{exo:b3:supramolecular:7}
An ITC experiment at \qty{298}{K} gives $K = \qty{2.0e5}{L.mol^{-1}}$ and $\Delta H^\circ = \qty{-40}{kJ/mol}$. Compute $\Delta G^\circ$, $T\Delta S^\circ$ and
$\Delta S^\circ$.
\end{exercise}

\begin{exercise}[$\star\star$]\label{exo:b3:supramolecular:8}
Explain the role of copper(I), and of its tetrahedral geometry, in Sauvage's
\omterm{def:b3:supramolecular:interlocked}{catenane} synthesis. How is the copper removed at the end?
\end{exercise}

\begin{exercise}[$\star\star$]\label{exo:b3:supramolecular:9}
The two signals of a \omterm{def:b3:supramolecular:interlocked}{rotaxane} shuttle, \qty{120}{Hz} apart, coalesce at \qty{300}{K} (data of the
exercise). Compute the exchange rate constant at coalescence and the \omterm{def:b3:rate-theories:activation}{Gibbs energy
of activation} of the shuttling.
\end{exercise}

\begin{exercise}[$\star\star\star$]\label{exo:b3:supramolecular:10}
A drug has an intrinsic solubility of \qty{0.10}{mmol/L} and forms a 1:1 complex with a
\omterm{def:b3:supramolecular:hosts}{cyclodextrin}, $K = \qty{500}{L.mol^{-1}}$ (data of the exercise). Show that the total solubility in a
solution of total \omterm{def:b3:supramolecular:hosts}{cyclodextrin} $C_t$ is $S_0 + KS_0C_t/(1 + KS_0)$, and compute it for $C_t = \qty{10}{mmol/L}$.
\end{exercise}

\begin{exercise}[$\star\star\star$]\label{exo:b3:supramolecular:11}
Show from the exact isotherm that when $KH_0 \gg 1$ the bound fraction equals the
number of equivalents up to one, then stays at one. Why does such a titration
give only a lower bound on $K$?
\end{exercise}

\begin{exercise}[$\star\star\star$]\label{exo:b3:supramolecular:12}
For the \omterm{def:b3:supramolecular:supramolecular}{host} of \cref{exo:b3:supramolecular:5}, assume the association is
\omterm{def:b3:rate-theories:diffusion}{diffusion-controlled}, $k_{\mathrm{on}} = \qty{1e9}{L.mol^{-1}.s^{-1}}$. Compute $k_{\mathrm{off}}$ and show that the exchange is fast on
the NMR time scale at \qty{400}{MHz}.
\end{exercise}

\section{Problem: Purple Benzene}

\begin{problem}[{Weekend problem --- purple benzene: potassium in a crown ether,
the extraction of permanganate into an organic solvent, a binding constant
from an NMR titration, and why a catalytic amount of crown suffices}]
\label{pb:b3:supramolecular:1}
% ledger: const:R, ghs:KMnO4, ghs:toluene, rion:K+VI, rion:Na+VI
Data of the problem. In methanol at \qty{298}{K}, $\log K = 6.1$ for $\ce{K+}$ and 4.4 for $\ce{Na+}$ with
18-crown-6. Extraction: an aqueous phase with \qty{0.10}{mol/L} $\ce{KCl}$ and \qty{1.0}{mmol/L} $\ce{KMnO4}$ is shaken
with an equal volume of toluene containing the crown L; the only extracted
species is the ion pair $\ce{[KL]+}\ce{MnO4-}$, with
$K_{\mathrm{ex}} = [\ce{[KL]+MnO4-}]_{\mathrm{org}}/([\ce{K+}]_{\mathrm{aq}}[\ce{MnO4-}]_{\mathrm{aq}}[\mathrm L]_{\mathrm{org}}) = \qty{1.0e5}{L^2.mol^{-2}}$, and the crown stays in the toluene. NMR titration of a
\omterm{def:b3:supramolecular:hosts}{crown ether} (\qty{2.0}{mmol/L}) with a sodium salt in deuterated methanol, one $\ce{CH2}$ signal
(ppm) against equivalents of $\ce{Na+}$:

\begin{center}\small
\begin{tabular}{l*{11}{c}}
\hline
equiv. & 0 & 0.25 & 0.5 & 0.75 & 1.0 & 1.5 & 2.0 & 3.0 & 4.0 & 6.0 & 8.0\\
$\delta$ & 3.560 & 3.590 & 3.612 & 3.635 & 3.649 & 3.674 & 3.688 & 3.697 & 3.704 & 3.708 & 3.714\\
\hline
\end{tabular}
\end{center}
% table: binding-titration-c

\medskip\noindent\textbf{Part I --- The complex.}
\begin{enumerate}
  \item Draw 18-crown-6 and its potassium complex.
  \item Compute $\Delta G^\circ$ of binding of $\ce{K+}$ and of $\ce{Na+}$.
  \item Compute the $\ce{K+}/\ce{Na+}$ selectivity.
  \item With the six-coordinate radii $r(\ce{K+}) = \qty{138}{pm}$ and $r(\ce{Na+}) = \qty{102}{pm}$, explain the selectivity.
  \item Why is the complex soluble in toluene when $\ce{KCl}$ is not?
  \item Why does the crown bind less strongly in water than in methanol?
\end{enumerate}

\medskip\noindent\textbf{Part II --- Extraction.}
\begin{enumerate}[resume]
  \item Write the extraction equilibrium.
  \item With $y$ the extracted concentration, write the equation for $y$ when
        $[\mathrm L]_0 = \qty{1.0}{mmol/L}$.
  \item Solve it, and give the fraction of permanganate extracted.
  \item Compute the fraction with $[\mathrm L]_0 = \qty{10}{mmol/L}$.
  \item Compute it with $[\mathrm L]_0 = \qty{0.10}{mmol/L}$.
  \item Why is the toluene purple, and why is the extracted permanganate a stronger
        oxidant than in water?
\end{enumerate}

\medskip\noindent\textbf{Part III --- An NMR titration.}
\begin{enumerate}[resume]
  \item Why does one averaged signal move instead of two signals?
  \item Write the model for $\delta$ against the \omterm{def:b3:supramolecular:supramolecular}{guest} concentration.
  \item Fit $K$ and $\delta_{\mathrm{bound}}$ by least squares; give $K$ with its standard error.
  \item Give the residual scatter and comment.
  \item Over what range of equivalents is the bound fraction between 20 and 80\,\%?
  \item Why would the same titration fail to measure $K$ for potassium, $\log K = 6.1$?
\end{enumerate}

\medskip\noindent\textbf{Part IV --- Phase-transfer catalysis.}
\begin{enumerate}[resume]
  \item An alkene in toluene is oxidised by the extracted permanganate. Where does
        the reaction take place?
  \item What happens to the crown after the reaction, and why does a catalytic
        amount suffice?
  \item Why is excess $\ce{KCl}$ in the water useful?
  \item What cheaper catalysts do the same job in industry?
  \item Why is toluene used today where Pedersen used benzene?
  \item State the result: the fraction of permanganate extracted at
        $[\mathrm L]_0 = \qty{1.0}{mmol/L}$.
\end{enumerate}
\end{problem}
